% Pursuing Stacks (cotes 134-2 to 134-8): archive cote 134-3, site page 31 (PDF page 32), typescript page 31.
% typescript. Draft compared: ps2.tex:1456-1532.
% AI-transcribed by Claude Opus 5.5 (claude-opus-5-5, adaptive thinking, effort high), Message Batches API, batch ps1 (msgbatch_01Rv5hCaNTtUt4JHx1yqoHPL), 2026-11-12; not yet checked by eye.
\scanpage{134-3/31}
\mspage{31}
one in the situation here. Let's call a ``modelizing category'', or sinply
a ``\uline{modelizer}'' ( ``modélisatrice'' in French), any category $M$, endowed with
a set $W\subset \mathrm{Fl}(M)$ (the weak equivalences), satisfying the obvious condition:

(M\unsure{\hand{o}}d) a) $W$ is the set of arrows made invertible by the localization
functor $M\to W^{-1}M$ \hand{(28')}, and
b) $W^{-1}M$ is equivalent to (Hot) ,
or equivalently, there exists a localization functor (*) (necessarily
unique up to unique isomorphism) such that $W$ be the set of arrows made
invertible by this functor.

Let $(M,W)$, $(M',W')$ be two modelizers, a functor \inserted{$F$:} $M\to M'$ is called
\uline{model-preserving}, or a \uline{morphism} between the modelizers, if it satisfies
either of the two following equivalent conditions:
(i) $F(W)\subset W'$\hand{,} hence a functor $F_{W,W'} : W^{-1}M\to W'^{-1}M'$, and the latter
is an equivalence.
(ii) The diagram
\diagram{\begin{tikzcd}[column sep=small]
M \arrow[rr, "F"] \arrow[dr] & & M' \arrow[dl] \\
& \text{(Hot)} &
\end{tikzcd}}
is commutative up to isomorphism (where the vertical arrows
are the \xout{\unsure{model}} ``type functors'' associated to $M$, $M'$ respectively\hand{)}

When dealing with a modelizer $(M,W)$, $W$ will be generally understood
so that we write simply \hand{$M$}. When $M$ is defined in terms of an algebraic
structure species, the task arises to find out whether (if any) there
exists a \uline{unique} $W\subset \mathrm{Fl}(M)$ turning $M$ into a modelizer, and if not so, if
we can however pinpoint one which is a more natural one, and which we
would call ``canonical''.

Here is a diagram including most of the modelizers and model-preserving
functors between these which we met so far\hand{,}\hand{(29)} \xout{with the} \hand{(}not included however
those connected with the theory of ``higher'' stacks and Gr-stacks, wh ich we
will have to elaborate upon later on\hand{)}:
\margin{$\Delta^\wedge$ (circled)}
\margin{$\square^\wedge$ (circled)}
\diagram{\begin{tikzcd}[row sep=small, column sep=small]
& & & \hand{$\hat{\Delta}$}=\text{(ss. sets)} \arrow[dr, "\hand{$\xi$}"] & \\
\text{(Ord)} \arrow[r, hook] & \text{(Preord)} \arrow[r, hook] & \text{(Cat)} \arrow[ur, "\hand{$\alpha$}"] \arrow[dr, "\hand{$\beta$}"'] & & \text{(Cat)} \\
& & & \hand{$\square$}=\text{(cub. sets)} \arrow[ur, "\hand{$\eta$}"'] &
\end{tikzcd}}
where the two last functors, with values in (Cat), are the two obvious
functors, obtained from
\hand{(**)} \hand{$i_A$}: $A^{\wedge}\to$ (Cat) , \over{a}{$F$} $\mapsto A_{/\over{a}{F}}$
by particularizing to $A=$\hand{$\Delta$} or \hand{$\square$}. As we noticed before, the four first
among these six functors admit left adjoints, but except for the first,
these adjoints are \uline{not} model preserving. The two last functors, and more
